\scanpage{21}
\mspage{18}
% archive sheet number: 20
\uline{Philosophie} Chaque fois qu'on a \deleted{deux \unsure{dérivateurs}} \inserted{un foncteur}
\[
\mathcal A\xrightarrow{u_0}\mathcal A'
\]
entre deux catégories qui peuvent s'interpréter \inserted{naturellement} comme les catégories fondamentales de deux dérivateurs $\mathbb D,\mathbb D'$, il se pose la question d'interpréter $u_0$ comme provenant d'un morphisme de dérivateurs
\[
\mathbb D\xrightarrow{u}\mathbb D'
\]
Celui-ci est, en pratique, tjrs déterminé canoniquement, à iso. canonique près. (Mais je doute pourtant que le foncteur ``restriction aux cat. fondamentales''
\[
\underline{\mathrm{Hom}}(\mathbb D,\mathbb D')\longrightarrow\underline{\mathrm{Hom}}(\mathcal A,\mathcal A')
\]
soit pleinement fidèle \ldots)

On va appliquer cette philosophie aux foncteurs $f^*$, $f_!$, $f_*$ (et éventuellement $f^!$, $f^?$), pour $f:I\longrightarrow J$, donc pour des foncteurs entre $\mathbb D(I)$ et $\mathbb D(J)$. \unsure{Cela|Ceci} nous amène donc à regarder $\mathbb D(I)$ et $\mathbb D(J)$ comme des cat. fondamentales de dérivateurs convenables.

\uline{Dérivateur }\uuline{induit} par un dérivateur $\mathbb D$, sur une catégorie d'indices $I\in\unsure{\mathrm{Dia}|\mathrm{Di}}$. C'est un dérivateur $\mathbb D_I$ défini par
\[
\mathbb D_I(J)=\mathbb D(I\times J)
\]
Le foncteur associé à $\mathbb D_I$, pour deux cat $J,J'\in\mathrm{Dia}$
